\section{Introdução}
\label{introducao}

% Parágrafos 1 e 2 (template): apresentação do tema dentro de um contexto.
A Engenharia de Software nasceu como disciplina na conferência da NATO de 1968 \citep{Naur1969NATO} e, desde então, responde a cada crise com novos processos, das metodologias ágeis ao DevOps. A chegada dos \textit{Large Language Models} (LLMs), treinados sobre grandes volumes de texto e código, abriu um novo ciclo dessa evolução. Ferramentas como o Cursor \citep{Cursor2022} e o Claude Code \citep{ClaudeCode2025} mostraram que modelos de linguagem interpretam intenções em linguagem natural e geram código para tarefas de implementação, depuração e testes. \cite{ES302024} chamam de Engenharia de Software 2.0 (ES 2.0) o conjunto de práticas em que o desenvolvedor usa esses assistentes como copilotos. Nesse paradigma, a IA acelera tarefas, mas o desenvolvedor conduz o fluxo de trabalho e decide.

Os copilotos têm limites conhecidos. \cite{ES302024} apontam a sobrecarga cognitiva do desenvolvedor, que permanece no centro do laço: decompõe o problema, instrui o assistente, avalia as sugestões e depura as falhas; apontam também a qualidade subótima do código gerado, e registram que edições em muitos arquivos continuam pouco confiáveis. Dois limites importam para este trabalho: o contexto, porque o desempenho dos modelos cai à medida que a janela cresce \citep{Levy2024SameTask, Liu2024LostInTheMiddle}, e a verificação, que nos copilotos fica inteira com o desenvolvedor. Em resposta, \cite{ES302024} propõem a Engenharia de Software 3.0 (ES 3.0), um paradigma centrado em intenção. Nele, o humano formula e refina intenções em conversa com agentes de IA, os \textit{AI Teammates}, que as transformam em software executável, e o ciclo termina quando o humano se dá por satisfeito. Sistemas como o MetaGPT \citep{MetaGPT2023} e o CodeAir \citep{CodeAir2025} materializam essa proposta: distribuem o trabalho entre agentes com papéis distintos, coordenados por um fluxo definido. A orquestração de papéis ataca os dois limites citados: uma memória de trabalho entre ciclos limita o contexto, e um revisor verifica o resultado antes da entrega.

Esses sistemas, porém, foram avaliados como um todo, com modelo, prompts e ferramentas diferentes dos sistemas com que foram comparados. Por isso, não se sabe o que a organização em papéis contribui por si. A maior parte da literatura compara sistemas inteiros, e falta um experimento em que só a topologia varie, com resolução verificada por testes escondidos e com medidas de custo, tempo, consistência, qualidade e integridade. Este trabalho delimita o tema por esse problema de pesquisa: a orquestração de papéis muda a correção de defeitos reais em relação a um agente único, com o modelo, as ferramentas, o prompt e o orçamento mantidos constantes?

O estudo fixou três hipóteses antes da coleta, aqui numeradas P1 a P3. A primeira (P1) afirma que a topologia muda a taxa de resolução das tarefas. Ela é neutra e bicaudal: não presume que o braço orquestrado seja melhor nem pior que o braço agente único. A segunda (P2) afirma que a topologia muda o custo em tokens nas tarefas resolvidas pelos dois braços. A terceira (P3) afirma que o efeito da topologia depende da faixa de tamanho da tarefa.

O objetivo geral deste trabalho é medir o efeito da topologia de agentes sobre a correção de defeitos reais, com o resto constante. Para alcançá-lo, definem-se os seguintes objetivos específicos:
\begin{itemize}
    \item construir o \textit{harness} Uranium, com oráculo de testes escondidos, detector de trapaça e caixa de areia;
    \item minerar e validar 13 tarefas de projetos Python abertos;
    \item calibrar e congelar os tetos de orçamento por run;
    \item executar as coletas pareadas por semente nos dois braços;
    \item replicar o experimento em três modelos; e
    \item analisar as sete métricas definidas antes da coleta.
\end{itemize}

A justificativa é prática e científica. Equipes que adotam agentes de código precisam decidir como organizá-los: um agente único com ferramentas ou um grafo de papéis especializados. Essa decisão depende de saber o que a topologia muda em resolução, em custo e no modo de falhar. A maior parte da literatura não isola essa variável, porque compara sistemas que diferem em várias dimensões ao mesmo tempo. Um experimento com o resto constante separa o efeito da topologia do efeito do modelo, do prompt e das ferramentas. O resultado orienta tanto a escolha de arquitetura quanto o desenho de estudos maiores.

A metodologia é o \textit{Design Science Research} (DSR), conforme \cite{Wieringa2009, Wieringa2014}, que combina a construção de um artefato com a sua avaliação empírica. O artefato é o Uranium, um \textit{harness} construído sobre LangGraph, com telemetria própria, oráculo de testes escondidos, detector de trapaça e caixa de areia. O Uranium é o meio da pesquisa, não o objeto. A avaliação é um experimento controlado, sem participantes humanos: 13 tarefas, 3 repetições pareadas por semente, 2 braços e 3 modelos, num total de 234 runs finais. Os três modelos vêm de três laboratórios (Alibaba, DeepSeek e Mistral) e são servidos por dois provedores (DeepInfra e Mistral). As sete métricas têm hierarquia fixada antes da coleta, com a resolução como desfecho primário.

O restante do texto está organizado em quatro seções. A Seção~\ref{referencial} apresenta os conceitos que fundamentam o estudo e os trabalhos correlatos. A Seção~\ref{desenvolvimento} descreve a caracterização da pesquisa, as tarefas, o Uranium, o projeto de experimentos e o ambiente. A Seção~\ref{resultados} apresenta, analisa e discute os resultados das três coletas. A Seção~\ref{conclusoes} reúne as conclusões, as contribuições e os trabalhos futuros.
% Os rótulos \label{resultados} e \label{conclusoes} precisam existir nos arquivos das seções 4 e 5.
